\documentclass[11pt]{article}

\usepackage{amsmath, amssymb}
\usepackage{booktabs}
\usepackage{geometry}
\usepackage{hyperref}
\usepackage{graphicx}

\geometry{margin=1in}

\title{STAT405 Project}
\author{Pierre Lardet}
\date{}

\begin{document}

\maketitle

\noindent Github: \url{https://github.com/PierreRL/STAT405_Project}

\section{Introduction}

In this project, I explore the causes of evacuation orders from Canadian wildfires. In particular, I seek to examine whether some populations (e.g.\ Indigenous populations) and/or regions are treated differently in terms of whether an evacuation order is issued.

\section{Data}

The primary dataset contains records of wildfire evacuation orders in Canada from 1980 to 2019 and the Canadian National Fire Database (NFDB) contains detailed wildfire characteristics of all wildfire regardless of whether an evacuation order was issued. I may also link these datasets with the Canadian Census to obtain demographic information about the affected communities.

\begin{itemize}
\item Evacuation order dataset (1980--2019): \\
\url{https://zenodo.org/records/5703323}

\item Canadian National Fire Database (NFDB), containing detailed wildfire characteristics: \\
\url{https://cwfis.cfs.nrcan.gc.ca/datamart/download/nfdbpnt}

\item Canadian Census data, which can provide demographic information about the affected communities. \url{https://www12.statcan.gc.ca/datasets/index-eng.cfm?Temporal=2021}
\end{itemize}

The heads for the evacuation order dataset and the NFDB dataset are shown below:

\begin{table}[h]
\centering
\resizebox{\textwidth}{!}{
\begin{tabular}{lllllllllllllllllll}
\toprule
EvacID & Year & Province & Location & Lat & Long & Population & PopulationNotes & FN\_Reserve & EvacSize & IssueDate & EndDate & EvacReason & Region & Group & FireSource & RepDOY & FireIgnit & FireSize \\
\midrule
1 & 1980 & MB & Clear Lake & 50.678605 & -99.911954 & NA & NA & no & Very Small (1--100) & 5/23/1980 & 5/27/1980 & Threat & SC & SC3 & NFDB.poly & 142 & Human & 22914.74086 \\
\bottomrule
\end{tabular}}
\caption{Example row from the wildfire evacuation order dataset.}
\end{table}

\begin{table}[h]
\centering
\resizebox{\textwidth}{!}{
\begin{tabular}{lllllllllllllllll}
\hline
FID & NFDBFIREID & SRC\_AGENCY & FIRE\_ID & FIRENAME & LATITUDE & LONGITUDE & YEAR & MONTH & DAY & REP\_DATE & OUT\_DATE & SIZE\_HA & CAUSE & FIRE\_TYPE & RESPONSE & MORE\_INFO \\
\hline
0 & BC-2023-2023-C23369 & BC & 2023-C23369 & C23369 & 51.8197 & -122.0729 & 2023 & 9 & 27 & 2023-09-27 & 2023-09-28 & 0.01 & H & Fire & FUL & NE of Alkali Creek \\
1 & BC-2023-2023-R50580 & BC & 2023-R50580 & KITSELAS &  &  &  &  &  &  &  &  &  &  &  &  \\
\hline
\end{tabular}
}
\caption{Example rows from the Canadian National Fire Database (CNFDB) fire point dataset.}
\end{table}

\newpage
\section{Model}

As a first model, I consider a Bayesian logistic regression:

\begin{align}
Y_i &\sim \mathrm{Bernoulli}(p_i) \\
p_i &= \operatorname{logistic}\!\left(\alpha + \beta^\top C_i + \gamma^\top F_i \right) \\
\alpha &\sim \mathcal{N}(0,\sigma_a^2) \\
\beta &\sim \mathcal{N}(0,\sigma_b^2) \\
\gamma &\sim \mathcal{N}(0,\sigma_c^2)
\end{align}
\begin{itemize}
\item $i$ indexes wildfires.
\item $Y_i$ is a binary variable indicating whether an evacuation order was issued.
\item $p_i$ is the probability that an evacuation order is issued.
\item $C_i$ represents characteristics of the affected community.
\item $F_i$ represents characteristics of the fire.
\item $\alpha$ is the intercept.
\item $\boldsymbol{\beta}$ and $\boldsymbol{\gamma}$ are coefficient vectors for community and fire characteristics.
\end{itemize}

The priors on $\alpha$, $\boldsymbol{\beta}$, and $\boldsymbol{\gamma}$ complete the Bayesian model. The variance parameters $\sigma_a^2$, $\sigma_b^2$, and $\sigma_c^2$ will be specified later.

Note that $C_i$ and $F_i$ may contain heterogeneous data types. Categorical variables may be converted to one-hot encodings, while continuous variables may be standardized prior to modelling.

More sophisticated models could include interaction terms between community and fire characteristics or hierarchical structure (e.g.\ random effects for region or province).

\section{Questions}

The primary scientific questions involve interpreting the parameters $\boldsymbol{\beta}$ and $\boldsymbol{\gamma}$.

\begin{itemize}
\item $\boldsymbol{\beta}$ describes how community characteristics affect the probability that an evacuation order is issued.
\item $\boldsymbol{\gamma}$ describes how wildfire characteristics influence the probability of evacuation.
\end{itemize}

A particular focus is whether certain community characteristics (e.g.\ Indigenous communities) are associated with a lower probability of evacuation orders, which could suggest systematic differences in evacuation decision-making.

\end{document}